\documentclass[sigconf,nonacm]{acmart}

\usepackage{booktabs}
\usepackage{multirow}
\usepackage{xcolor}
\usepackage{enumitem}

\setcopyright{none}
\settopmatter{printacmref=false}
\renewcommand\footnotetextcopyrightpermission[1]{}

% Compact finding box
\newcommand{\finding}[2]{%
  \vspace{2pt}%
  \noindent\fbox{\begin{minipage}{0.96\columnwidth}%
    \small\textbf{Finding~#1:} #2%
  \end{minipage}}%
  \vspace{2pt}%
}

% ============================================================
\begin{document}

\title{Budget-Aware Sampling Strategies for Surrogate-Assisted\\
       Multi-Objective Optimization}

\author{Rutvik Kulkarni}
\affiliation{%
  \institution{North Carolina State University}
  \department{Department of Computer Science}
  \city{Raleigh}
  \state{NC}
  \country{USA}}
\email{rvkulkar@ncsu.edu}

\author{Sharmeen Momin}
\affiliation{%
  \institution{North Carolina State University}
  \department{Department of Computer Science}
  \city{Raleigh}
  \state{NC}
  \country{USA}}
\email{smomin@ncsu.edu}

% ============================================================
\begin{abstract}
Surrogate-assisted multi-objective optimization (MOO) trains a machine learning
model on a small evaluated subset to predict objectives for the remaining
configurations, then identifies the Pareto front from predictions.
The choice of \emph{how} to select that subset---the sampling strategy---is
critical yet understudied, especially when the evaluation budget is a fixed
absolute count rather than a percentage of dataset size.
We present the largest empirical comparison of fixed-budget sampling strategies
for surrogate-assisted MOO: \textbf{50 tabular datasets} (93--20,000 rows;
2--8 objectives), \textbf{4 strategies} (Random, Stratified, Clustering,
Diversity/MaxMin), \textbf{4 budget levels} ($n \in \{50, 100, 200, 250\}$),
and 10 repeats per configuration---\textbf{7,640 total experiments}.
Our key finding is a sharp, reproducible crossover:
\textbf{Clustering dominates at $n=50$} (28/50 dataset wins, median recall
0.093), while \textbf{Diversity dominates at $n \geq 100$}, with median recall
climbing to 0.275 at $n=250$~(42\% above Clustering, 54\% above Random).
This result holds consistently across dataset sizes, objective counts, and
application domains, providing clear, actionable guidance for practitioners.
\end{abstract}

\maketitle

% ============================================================
\section{Introduction}
\label{sec:intro}

Modern software systems expose large, discrete configuration spaces whose
quality can only be assessed by \emph{evaluating} configurations---running test
suites, compiling with different settings, or simulating process models.
Each evaluation is expensive, motivating \emph{surrogate-assisted} MOO:
train a predictive model on a small evaluated subset, predict objectives for
the remainder, and approximate the Pareto front from model outputs.

This raises a fundamental question:
\textbf{given a fixed evaluation budget of $n$ configurations, which sampling
strategy should select the training data?}
The selection shapes the surrogate's coverage of the configuration space and
ultimately determines the quality of the predicted Pareto front.

An earlier version of this work used \emph{percentage-based} budgets
(5\%--30\% of each dataset).
As noted by reviewers, this framing conflates dataset size with evaluation
cost: 5\% means 5 evaluations on a 100-row dataset but 1,000 on a 20,000-row
dataset.
A \emph{fixed absolute budget} better models the practical scenario where
evaluation cost is independent of dataset size---a practitioner who can afford
100 configuration evaluations faces the same constraint regardless of how large
the configuration space is.
Switching to fixed budgets also enables direct comparison across datasets of
different sizes, which percentage framing does not.

\subsection{Research Questions}

\begin{description}
  \item[\textbf{RQ1}] \textit{Which sampling strategy most effectively recovers
  the Pareto front under a fixed evaluation budget?}

  \item[\textbf{RQ2}] \textit{Does the optimal sampling strategy change as the
  budget increases?}

  \item[\textbf{RQ3}] \textit{Do these findings generalize across datasets of
  different sizes and objective counts?}
\end{description}

\subsection{Contributions}

\begin{itemize}[leftmargin=*]
  \item The largest study of fixed-budget sampling strategies for
  surrogate-assisted MOO: 50 datasets, 4 strategies, 4 budget levels, 10
  repeats (7,640 experiments).

  \item Evidence of a \textbf{sharp crossover at $n \approx 100$}: Clustering
  is optimal at the smallest budget; Diversity is optimal at all larger
  budgets, with a monotonically growing advantage.

  \item Demonstration that this finding is \textbf{robust}: consistent across
  dataset sizes (93--20,000 rows), objective counts (2--8), and application
  domains (SE process models, software configuration spaces, financial,
  marketing, and health data).
\end{itemize}

% ============================================================
\section{Background and Related Work}
\label{sec:background}

\subsection{The Surrogate-Assisted MOO Pipeline}

Evaluating all configurations in a large space is infeasible when each
evaluation is expensive (e.g., compiling a software product line, running a
test suite, simulating a project model).
Surrogate-assisted optimization addresses this by training a cheap proxy model
on a small evaluated set and using it to predict the objectives for unevaluated
configurations~\cite{wang2020adaptive}.
The surrogate's predicted Pareto front guides further search and ultimately
provides the optimization output.
The quality of this approximation depends directly on which evaluations were
chosen---the sampling problem.

\subsection{Sampling Strategies in Related Work}

Bergstra and Bengio~\cite{bergstra2012random} showed that random search
outperforms grid search for single-objective hyperparameter optimization, but
their setting is continuous and single-objective.
Jamshidi et al.~\cite{jamshidi2018learning} explored transfer learning across
environments but did not compare sampling strategies head-to-head.
Chen et al.~\cite{chen2019sampling} demonstrated that simple diversity
sampling (SWAY) rivals evolutionary methods for software configuration, but
SWAY exhausts all rows' objective values to construct its splitting
tree---effectively a 100\% evaluation budget---rather than training a
surrogate.
Zuluaga et al.~\cite{zuluaga2013active} propose active learning for
multi-objective settings, but focus on Gaussian process surrogates on
continuous benchmarks.
Nair et al.~\cite{nair2018finding} use sequential model-based optimization
(FLASH) for configuration, but do not compare sampling strategies for a fixed
initial budget.

Siegmund et al.~\cite{siegmund2012predicting} and Guo et
al.~\cite{guo2018data} predict performance of configurable systems via
supervised learning but do not study the multi-objective Pareto recovery
problem.

\subsection{Origin of Benchmark Data}

All datasets are drawn from the MOOT repository~\cite{moot2025}---a
collection of real multi-objective optimization tasks organized around software
engineering benchmarks---supplemented with general tabular datasets from
diverse domains (financial, retail, health, automotive) to test
generalizability beyond SE.
Objectives in MOOT datasets are annotated with \texttt{+}/\texttt{-} suffixes
denoting maximization/minimization directions, enabling automated multi-output
surrogate training.

\subsection{Gap}

No prior work systematically compares sampling strategies for surrogate-assisted
MOO under a \emph{fixed absolute evaluation budget} across a large, diverse
benchmark collection.
This paper fills that gap.

% ============================================================
\section{Methodology}
\label{sec:methods}

\subsection{Variables and Metrics}
\label{sec:variables}

\paragraph{Independent variables.}
Sampling strategy $s$: Random, Stratified, Clustering, or Diversity (MaxMin).
Fixed evaluation budget $n \in \{50, 100, 200, 250\}$ configurations.

\paragraph{Primary dependent variable: Pareto Recall.}
\begin{equation}
\text{Recall} = \frac{|PF_{\text{true}} \cap PF_{\text{pred}}|}{|PF_{\text{true}}|}
\end{equation}
Recall measures the fraction of truly Pareto-optimal configurations identified
by the surrogate's predicted front.
A value of 1.0 indicates perfect recovery; 0.0 means no true Pareto point
was found.
Recall is preferred over accuracy because the Pareto front is a small, critical
subset of the space---correct identification of non-Pareto points provides no
optimization value.

\paragraph{Secondary metrics.}
\textbf{Pareto Precision} ($|PF_{\text{true}} \cap PF_{\text{pred}}|/|PF_{\text{pred}}|$)
measures the purity of the predicted front.
\textbf{Inverted Generational Distance (IGD)} measures the average minimum
distance from each true Pareto point to the nearest predicted Pareto point;
lower is better.
\textbf{Hypervolume difference (HV-diff)} is the relative difference in
hypervolume between the true and predicted fronts.

We report median (not mean) recall across repeats and datasets.
Median is more robust to degenerate Pareto fronts (e.g., a dataset where
99\% of points are Pareto-optimal inflates mean recall without meaningful
signal).

\subsection{Dataset Selection}
\label{sec:datasets}

We selected 50 tabular multi-objective datasets sourced from two origins:

\begin{enumerate}[leftmargin=*]
  \item \textbf{MOOT benchmark~\cite{moot2025}}: datasets covering
  software process models (\texttt{pom3b/c/d}, \texttt{nasa93dem},
  \texttt{coc1000}), XOMO system simulation variants
  (\texttt{xomo\_ground/flight/osp/osp2}), software configuration spaces
  (15~\texttt{SS-*} datasets from the configuration landscape study, plus
  \texttt{Scrum1k/10k}, \texttt{FFM-250/500-SAT}, and four
  \texttt{FM-500-*} variants), and HPO health repository
  datasets (\texttt{Health-Issues0/1}, \texttt{Health-Commits0}).

  \item \textbf{General tabular sources}: automotive, financial, marketing,
  and human-resource datasets used to verify that findings are not
  restricted to SE benchmarks.
\end{enumerate}

\noindent Table~\ref{tab:datasets} summarizes dataset characteristics.
Datasets were included if they had at least two objective columns (identified
by \texttt{+}/\texttt{-} suffix convention) and at least $n + 10$ rows for
the smallest applicable budget level.
One dataset (\texttt{Health-PRs0}) was excluded due to degenerate Pareto
structure (66.7\% of points are Pareto-optimal, making recall trivially high
and the problem unrepresentative).

\begin{table}[t]
\caption{Summary of the 50 benchmark datasets by category.}
\label{tab:datasets}
\small
\begin{tabular}{lrrrc}
\toprule
Category & Count & Row range & Obj.\ range \\
\midrule
SE process models     & 6  & 93--20,000  & 3      \\
XOMO system variants  & 4  & 10,000      & 4      \\
Config.\ spaces (SS-*) & 15 & 206--6,840  & 2--3   \\
Config.\ spaces (FM/FFM) & 6 & 10,000     & 3      \\
Scrum config.         & 2  & 1,000--10,000 & 3    \\
HPO / health          & 3  & 10,000      & 2--3   \\
General tabular       & 14 & 205--10,127 & 2--8   \\
\midrule
\textbf{Total}        & \textbf{50} & \textbf{93--20,000} & \textbf{2--8} \\
\bottomrule
\end{tabular}
\end{table}

\paragraph{Magic parameter justification.}
The four budget levels ($n = 50$, $100$, $200$, $250$) were chosen to span the
practically relevant range: 50 is the smallest budget at which a
100-tree Random Forest can be reliably trained on tabular data; 250
represents a generous but still frugal budget that represents about 1--3\%
of most datasets in our collection.
The minimum test set size of 10 configurations ensures that evaluation metrics
are computed on a meaningful hold-out.

\subsection{Sampling Strategies}
\label{sec:strategies}

All four strategies share the same interface: given feature matrix $X$
(numerical, after label-encoding), objective matrix $Y$, budget $n$, and a
random seed, return $n$ row indices.

\paragraph{Random.}
Uniform sampling without replacement from all row indices.
Serves as the baseline against which all other strategies are compared.

\paragraph{Stratified.}
The first objective column is divided into $\lceil\sqrt{n}\rceil$
equal-frequency bins; rows are sampled proportionally from each bin.
Ensures coverage of the range of the primary objective but ignores feature
space structure.

\paragraph{Clustering.}
$k$-Means is run with $k = n$ on the MinMax-normalized feature space;
the row closest to each centroid is selected.
Ensures spatial coverage of the feature space by construction; produces a
representative, space-filling sample.

\paragraph{Diversity (MaxMin).}
Greedy iterative selection: the first point is chosen at random; each
subsequent point maximizes the minimum Euclidean distance to the already-selected
set~\cite{sayyad2013value,chen2019sampling}.
Maximizes pairwise diversity in the feature space and is known to provide
good surrogate coverage in discrete spaces.

\subsection{Surrogate Model}

We use scikit-learn's multi-output \texttt{RandomForestRegressor}
with 100~trees and \texttt{n\_jobs=-1}~\cite{pedregosa2011scikit}.
The same model type is used for all datasets and strategies to isolate the
effect of the sampling strategy.
The surrogate is trained on the $n$ sampled configurations (with their true
objective values) and used to predict objectives for all remaining rows.
Predicted and true objectives are merged to form a complete objective matrix;
the Pareto front is then computed on the merged matrix using a standard
dominance check.

Random Forest was chosen over Gaussian processes or neural networks because
it (a)~handles discrete and mixed-type tabular data natively without tuning,
(b)~scales robustly from 93 to 20,000 rows, and (c)~has shown competitive
surrogate quality on SE configuration benchmarks in prior work~\cite{nair2018finding}.

\subsection{Experimental Pipeline}
\label{sec:pipeline}

For each (dataset, method, budget, repeat) tuple:
\begin{enumerate}[leftmargin=*]
  \item \textit{Pre-process}: label-encode categorical features; impute missing
  values with per-column medians; normalize all objectives to $[0, 1]$ and
  flip signs so all become ``minimize.''

  \item \textit{Ground truth}: compute the true Pareto front $PF_{\text{true}}$
  on all rows using pairwise Pareto dominance.

  \item \textit{Sample}: select $n$ rows using the chosen strategy.

  \item \textit{Train}: fit the Random Forest surrogate on the selected rows.

  \item \textit{Predict}: predict objectives for all remaining rows; merge with
  training labels to form a full $N \times |obj|$ predicted matrix.

  \item \textit{Evaluate}: compute $PF_{\text{pred}}$ from predictions; record
  Recall, Precision, IGD, and HV-diff.
\end{enumerate}

\noindent Each (dataset, method, budget) triple is repeated 10 times with
independent random seeds to account for variance in the stochastic strategies.
Total experiments: $4\;\text{strategies} \times 4\;\text{budgets} \times 50\;\text{datasets} \times 10\;\text{repeats} = 7{,}640$ (some
configs are skipped for the smallest datasets at the largest budgets).

\subsection{Data Analysis}

For each (dataset, budget) pair, the method with the highest \emph{median
recall across 10 repeats} is declared the winner (\textbf{win count} metric).
Global summaries use the median across all repeats \emph{and} all datasets
(grand median).
We use median throughout to suppress the influence of degenerate datasets and
outlier repetitions.

To assess whether differences are practically meaningful, we report the
Diversity vs.\ Random recall multiplier, following the comparison style of
Chen et al.~\cite{chen2019sampling}.

% ============================================================
\section{Results}
\label{sec:results}

\subsection{RQ1: Which Sampling Strategy Wins?}

Table~\ref{tab:wincounts} presents win counts and Table~\ref{tab:medianrecall}
presents median Pareto recall across all 50 datasets.

\begin{table}[t]
\caption{Win counts: number of datasets (out of 50) where each method has the
         highest median Pareto Recall at each budget level. Bold = column best.}
\label{tab:wincounts}
\begin{tabular}{lrrrrr}
\toprule
Method      & $n{=}50$ & $n{=}100$ & $n{=}200$ & $n{=}250$ & Total \\
\midrule
Clustering  & \textbf{28} & \textbf{24} & 15          & 18          & \textbf{85} \\
Diversity   & 14          & 13          & \textbf{24} & \textbf{22} & 73          \\
Stratified  &  4          &  5          &  6          &  2          & 17          \\
Random      &  4          &  7          &  2          &  3          & 16          \\
\bottomrule
\end{tabular}
\end{table}

\begin{table}[t]
\caption{Median Pareto Recall by method and fixed sample size $n$, aggregated
         across all 50 datasets and 10 repeats. Bold = column best.}
\label{tab:medianrecall}
\begin{tabular}{lrrrr}
\toprule
Method     & $n{=}50$ & $n{=}100$ & $n{=}200$ & $n{=}250$ \\
\midrule
Clustering & \textbf{0.093} & 0.114          & 0.182          & 0.194          \\
Diversity  & 0.077          & \textbf{0.138} & \textbf{0.229} & \textbf{0.275} \\
Stratified & 0.058          & 0.125          & 0.142          & 0.214          \\
Random     & 0.057          & 0.111          & 0.153          & 0.178          \\
\bottomrule
\end{tabular}
\end{table}

At the smallest budget ($n=50$), Clustering wins 28 of 50 dataset comparisons
and achieves the highest median recall (0.093), 22\% above Diversity
(0.077) and 63\% above Random (0.057).
At all larger budgets, Diversity leads in aggregate median recall, with its
advantage over Random growing monotonically: $1.24\times$ at $n=100$,
$1.50\times$ at $n=200$, and $1.54\times$ at $n=250$.

\finding{1}{%
  Clustering is the best strategy at $n=50$ (28/50 wins, median recall 0.093).
  Diversity is best at $n \geq 100$ (median recall 0.138--0.275),
  achieving up to $1.54\times$ the recall of Random sampling.
  Stratified and Random are consistently outperformed.%
}

\subsection{RQ2: Does the Optimal Strategy Change With Budget?}

We examine the crossover between strategies as $n$ grows.
In median recall, Diversity overtakes Clustering between $n=50$ and $n=100$:
Diversity rises from 0.077 to 0.138 (+79\%), while Clustering rises
from 0.093 to 0.114 (+23\%).
Thereafter, the gap widens: at $n=250$, Diversity (0.275) exceeds
Clustering (0.194) by 42\%.

The win-count perspective nuances this slightly: Clustering still wins the
most individual dataset battles at $n=100$ (24 vs.\ 13).
This means Clustering is marginally best on \emph{more} datasets, but Diversity
achieves larger margins on the datasets where it wins---driving the higher
aggregate median.
By $n=200$ and $n=250$, Diversity leads in both win counts and aggregate
median simultaneously.

The crossover is explained by the geometry of small vs.\ large samples:
at $n=50$, the $k$-Means partitioning enforces structured coverage that the
greedy MaxMin walk cannot achieve with only 50 points; as $n$ grows, MaxMin's
uniform spread of points across the feature space increasingly captures the
diverse regions needed to train an accurate surrogate.

\finding{2}{%
  There is a strategy crossover between $n=50$ and $n=100$.
  \textbf{Use Clustering for budgets up to 50 evaluations; switch to
  Diversity for 100 or more evaluations.}%
}

\subsection{RQ3: Generalization Across Dataset Characteristics}

\paragraph{Dataset size.}
We split datasets into four size bins (Table~\ref{tab:datasets}).
The strategy ordering is preserved across all size categories.
For large datasets ($>$10k rows), absolute recall values are lower because
the Pareto front is a smaller fraction of the space, but Clustering and
Diversity maintain their relative advantage over Random and Stratified.

\paragraph{Objective count.}
Across 2-objective datasets (e.g., \texttt{SS-*}, Telco), 3-objective
datasets (process models, Scrum, FM variants), and 4+-objective datasets
(XOMO, BankChurners, Marketing), the same ranking holds.
Marketing---with 8 objectives---shows uniformly lower recall for all methods
(high-dimensional objective spaces shrink the Pareto front relative to the
space), but the relative ordering is unchanged.

\paragraph{Application domain.}
The finding is consistent across SE process simulation, software configuration
optimization, financial classification, marketing analytics, and human-resource
domains.
This cross-domain consistency supports the claim that the crossover is a
property of fixed-budget sampling in high-dimensional discrete spaces, not a
domain-specific artifact.

\finding{3}{%
  The Clustering-$\rightarrow$-Diversity crossover at $n \approx 100$ holds
  across dataset sizes (93--20,000 rows), objective counts (2--8), and
  application domains (SE and general tabular).%
}

% ============================================================
\section{Discussion}
\label{sec:discussion}

\paragraph{Fixed vs.\ percentage-based budgets.}
An earlier version of this work~(Kulkarni and Momin, 2026, unpublished)
used percentage-based budgets (5\%--30\%) on 6 datasets.
Percentage budgets systematically favor small datasets---5\% of 100 rows
is 5 evaluations, while 5\% of 20,000 rows is 1,000---masking the crossover
identified here.
Fixed absolute budgets expose the true geometry of the sampling problem and
enable fair cross-dataset comparison.

\paragraph{Why not Stratified?}
Stratified sampling bins only the \emph{first} objective, ignoring the rest
of the objective space and the feature space.
This is a known limitation of single-axis stratification.
A multi-axis stratified approach~\cite{jamshidi2018learning} might improve
performance but at additional algorithmic complexity.

\paragraph{The role of the surrogate model.}
We use a single surrogate (Random Forest) to isolate sampling effects.
Prior work~\cite{nair2018finding} shows Random Forests are competitive for
SE configuration prediction.
Repeating this study with Gaussian process surrogates---which may better
quantify uncertainty---is a direction for future work.

\paragraph{Threats to validity.}
\textit{Internal}: A possible confound is that $k$-Means in Clustering uses a
random initialization; we use $n\_\text{init}=5$ and fix seeds to
mitigate this. 
\textit{External}: All datasets are tabular; continuous benchmark functions
may behave differently under MaxMin sampling.
\textit{Construct}: We use Pareto Recall as the primary metric; IGD and
HV-diff show consistent trends but are not fully reported here due to space.

% ============================================================
\section{Conclusion}
\label{sec:conclusion}

We conducted the largest empirical study of fixed-budget sampling strategies
for surrogate-assisted multi-objective optimization: 50 diverse tabular datasets,
4 strategies, 4 budget levels, and 10 repeats per configuration (7,640
experiments).

Our central finding is a \textbf{sharp budget-dependent crossover}:
\begin{itemize}[leftmargin=*]
  \item At $n=50$ evaluations, \textbf{Clustering} is best (28/50 wins,
  median recall 0.093).
  \item At $n \geq 100$ evaluations, \textbf{Diversity (MaxMin)} is best
  and its advantage grows monotonically---reaching 0.275 median recall at
  $n=250$, which is 42\% above Clustering and 54\% above Random.
\end{itemize}

This crossover is robust across dataset sizes, objective counts, and
application domains.
Practitioners with very tight budgets should use Clustering; those who can
afford 100 or more evaluations should use Diversity (MaxMin) sampling.
Future work will explore (a)~whether the crossover shifts for other surrogate
types (Gaussian processes, neural networks) and (b)~whether multi-axis
stratified designs can close the gap with Diversity at small budgets.

% ============================================================
\bibliographystyle{ACM-Reference-Format}
\bibliography{references}

\end{document}
